\section{Introduction}
\label{sec:intro}

A \emph{hammer} translates a goal of a proof assistant, together with a
selection of library facts, into the input logic of automated theorem
provers (ATPs), and reconstructs any proof the ATPs find
\cite{BlanchetteEtAl2016qed}.  For proof assistants based on the Calculus
of Inductive Constructions (CIC) --- Rocq (formerly Coq) and its relatives
--- the reference tool is CoqHammer \cite{CzajkaKaliszyk2018,CoqHammer},
whose target is untyped classical first-order logic (FOL).  The efficiency
of such a translation hinges on it being \emph{shallow}
\cite{Czajka2018}: source propositions must become first-order
\emph{formulas} whose connectives, quantifiers and equalities are the
native first-order ones, so that clausification, superposition and term
indexing do real logical work on the translated problem, instead of
simulating the object logic through an inhabitation predicate one bridge
axiom at a time.

A companion paper \cite{Extraction2026} carries this programme out for a
monomorphic fragment of CIC with \emph{non-indexed} inductive datatypes:
the translation of a definition is factored into \emph{program extraction}
(the Letouzey-style erasure \cite{Letouzey2003,Letouzey2004} of the body,
compiled to per-constructor rewrite equations) and \emph{specification
extraction} (a realizability reading of the type as a first-order formula
about the translated constant), and both components are proved sound
against a single first-order model: a proof-irrelevant, classical term
model over the erasure calculus whose domain deliberately contains
ill-typed ``junk'' elements.  The model decides, uniformly, which guards
are semantically mandatory and which may be dropped.

That development stops at a well-marked boundary: its datatypes have no
indices.  \emph{Indexed inductive families} --- length-indexed vectors
$\vect\,A : \nat \to \Set$, finite sets $\Fin : \nat \to \Set$, bounded
numbers, the ssreflect view type $\mathsf{reflect}\ P : \bool \to \Set$,
and the user-declared equality-like and $\mathsf{le}$-like predicates ---
are the workhorses of dependently typed libraries, and they are exactly
the point where CIC's typing discipline diverges most sharply from
first-order logic: the type of a constructor \emph{constrains} the indices
at which it may appear, case analysis must \emph{unify} the indices of the
scrutinee's type with those of each constructor, and eliminations of
propositional families into computational content are governed by the
singleton criterion, Streicher's K, and their refinements
\cite{GoguenEtAl2006,CockxEtAl2014,GilbertEtAl2019}.  The present paper
extends the companion's calculus, translation and soundness proof to
indexed families, and settles the design questions specific to them.

\subsection{The idea: fording}
\label{sec:intro-idea}

The mechanism is due to McBride \cite[\S 3.5]{McBride1999}, who named it
after Henry Ford's ``any colour\,\dots\,so long as it is black'': a
constructor of an indexed family,
\[
  \wcns{c} : \forall \tup{a} : \tup{A}.\ I\,\tup{t}(\tup{a}),
\]
is equivalent to a constructor of a \emph{non-indexed} family that accepts
arbitrary indices $\tup{u}$ but carries equational premises restoring the
lost specificity:
\[
  \wcns{c}^{\mathrm{Ford}} :
  \forall \tup{u}.\ \forall \tup{a} : \tup{A}.\
  (\tup{u} = \tup{t}(\tup{a})) \to I^{\mathrm{Ford}}\,\tup{u}.
\]
Index constraints thereby become \emph{propositional payload}, on the same
footing as the refinement payloads $P\,x$ of subset types
$\refty{x{:}A}{P\,x}$, which the companion translation already expands
shallowly at every occurrence.  Every system that elaborates dependent
pattern matching performs this reduction internally --- Coq's
\texttt{inversion} \cite{CornesTerrasse1995,Monin2010}, the
Goguen--McBride--McKinna translation of pattern matching to eliminators
(``the equations on the indices are exactly those we must unify''
\cite{GoguenEtAl2006}), Agda's unifier \cite{CockxEtAl2014,CockxDevriese2018},
and the Equations package \cite{SozeauMangin2019}.  The present paper
performs it in a translation to untyped first-order logic, \emph{per
occurrence and per constructor}, and proves the result sound.

Two classical obstacles make fording delicate inside type theory, and both
dissolve in an untyped first-order target.  First, telescopic index
equations relate terms whose \emph{types} are only equal after earlier
equations have been substituted; stating them requires heterogeneous
(``John Major'') equality, whose eliminator is exactly Streicher's K
\cite{McBride1999,CockxEtAl2014}.  Untyped first-order equality states
them directly.  Second, defined functions in index positions ---
McBride's \emph{green slime} \cite{McBride2014} --- make elaboration-time
unification stuck: $0 \feq \mathsf{plus}(m,n)$ is neither a positive nor
a negative unification success because $\mathsf{plus}$ is not a rigid
symbol.  A saturation prover simply \emph{keeps} the equation and reasons
with the defining equations of $\mathsf{plus}$; specialization by
unification becomes ordinary superposition, and branches that a
type-theoretic elaborator must leave stuck can be closed by proof.  The
translation of this paper accordingly never solves index constraints; it
only states them.

What does \emph{not} dissolve by itself is soundness.  The literature on
translations of inductive types into untyped or weakly typed logics is a
catalogue of documented accidents precisely at the three points this paper
must get right: unguarded constructor-exhaustiveness bounds the
cardinality of the whole untyped universe (Meng and Paulson's two-element
universe \cite[\S 2.3]{MengPaulson2008}; the monotonicity analysis of
\cite{BlanchetteEtAl2016} identifies exhaustiveness axioms as exactly the
sentences that can never lose their guards); encoding a \emph{type-indexed}
equality by the native untyped equality of the target lets an equation
established at one index leak to another (the \fstar{} unsoundness
\#1542 \cite{FStar1542}, where pointwise function equality proved at
domain $\mathsf{nat}$ escaped to domain $\mathsf{int}$, yielding
$\mathsf{False}$); and axioms that are individually plausible ---
selector totality, rank/height orderings, box/unbox roundtrips --- can
interact inconsistently on degenerate or empty instances (the Verus
unsoundness \#1366 \cite{Verus1366}).  Our response to all three is the
companion's discipline, extended: \emph{every} emitted axiom is proved
valid in one fixed model, constructed once, so no combination of emitted
axioms can be inconsistent (\cref{cor:consistency}), and the model
arbitrates which guards are load-bearing.  Where a design variant is
unsound, we do not merely avoid it: we prove its unsoundness inside the
framework, as a specification for implementations
(\cref{sec:formulas}).

\subsection{Contributions}
\label{sec:intro-contrib}

This paper is self-contained: all definitions and proofs are given in
full, including the parts shared with the companion paper
\cite{Extraction2026}, which are re-developed here in their generalized
form.  Remaining citations to \cite{Extraction2026} in proofs indicate
provenance; none replaces an argument.  Concretely:

\begin{enumerate}[leftmargin=2em]
\item We define $\CICi$ (\cref{sec:calculus}), a stratified, monomorphic,
  first-order fragment of CIC whose datatypes are \emph{indexed inductive
  families} with dependent constructor telescopes and propositional
  constructor premises, subsuming the non-indexed datatypes, refinement
  types and indexed inductive predicates of the companion fragment.  Case
  analysis on a family is \emph{forded}: each branch binds, besides the
  constructor's arguments and premises, proofs of the index equations
  $\tup{u} = \tup{t}_i(\tup{y}_i)$, so that dependently typed programming
  idioms (impossible branches, index-directed results) are expressible
  with first-order motives, without large elimination.  The calculus also
  has elimination of \emph{index-determined} propositional singletons ---
  the image of CIC's singleton elimination for equality-like families,
  with McBride-style forcing in place of CIC's parameter side condition
  --- and a constructor-discrimination proof former.
\item We define the erasure of $\CICi$ into the untyped calculus $\lbox$
  and prove the two syntactic pillars of extraction --- substitution
  commutation and simulation of reduction --- together with confluence of
  $\lbox$ (\cref{sec:lambdabox}, \cref{sec:confluence}).  Erasure is
  index-blind; in particular proofs occurring \emph{inside index terms}
  vanish, which is the translation-level form of canonicalizing erased
  data to a single token \cite{Atkey2018,Werner2008}.
\item We define the translation (\cref{sec:translation}).  Guard
  predicates become $(1{+}m)$-ary --- the guard of an occurrence
  $I(\tup{u})$ is the atom $\Ihat{I}(w, \compf(\erase{\tup{u}}))$, carrying
  the translated indices --- and the structural axioms of a family are the
  first-order image of its ford: constructor typing with index conclusions,
  and a guarded \emph{inversion} axiom whose disjuncts carry constructor
  payloads \emph{and index equations}.  Pattern-matching definitions
  compile, unchanged, to unguarded per-constructor equations; the
  compilation layer itself is the companion's, with a repair to its
  treatment of recursive closures (\cref{rem:rec-closure}) that is to
  be mirrored there.  We prove
  the forced-argument simplification (index equations solvable for a bare
  variable are redundant) as a first-order equivalence
  (\cref{lem:n-one}).
\item We construct the model (\cref{sec:semantics}): the domain is again
  the set of closed $\lbox$ terms modulo conversion, junk included; an
  indexed family denotes a set of (index-tuple, element) pairs built in
  stages.  The key new lemma is \emph{index reflection}
  (\cref{lem:index-reflection}): every member of the interpretation
  decomposes uniquely as a constructor form whose index tuple equals the
  denotation of the constructor's index terms --- the semantic content of
  fording, and the point where the model validates what uniqueness of
  identity proofs validates in type theory.  Ranks become element ranks,
  which index reflection makes well defined.
\item We prove soundness (\cref{sec:soundness}): every axiom of the
  emitted theory holds in the model (\cref{thm:env-validity}), first-order
  provability of the translated goal implies truth of the goal in the
  proof-irrelevant classical semantics (\cref{thm:soundness}), the theory
  is consistent (\cref{cor:consistency}), and no refutable goal becomes
  provable (\cref{cor:no-refutable}).  The fundamental lemma
  (\cref{lem:fundamental}) is proved in full; its new cases --- forded
  case analysis, singleton elimination, fixpoints over families --- are
  exactly where index reflection is spent.  On premise-free families with
  first-order arguments the model is standard: $\semr{\vect}{}$ at index
  $n$ is in bijection with length-$n$ lists (\cref{prop:standard}).
\item We settle the indexed-specific design questions
  (\cref{sec:formulas}), each as a theorem about the model or the theory:
  \begin{itemize}
  \item the inversion axiom's guard is mandatory in the strongest
    sense: for \emph{every} indexed family, the unguarded axiom
    together with the family's own injectivity and discrimination
    axioms collapses the universe to a point, so it is inconsistent
    the moment any two-constructor datatype is declared; for families
    that specialise an index to a constructor form ($\vect$, $\Fin$),
    discrimination alone detonates it
    (\cref{prop:unguarded-inconsistent}, exact delimitation in
    \cref{rem:unguarded-general});
  \item the index arguments of guard predicates are mandatory: the
    index-forgetting variant proves a goal that is false in the source
    semantics --- the precise analogue of \fstar{} \#1542
    (\cref{prop:index-dropping});
  \item per-constructor match equations remain valid without any guard or
    index premise, \emph{including} equations for statically impossible
    branches; the index-premised variant is strictly weaker and both are
    sound, so the choice is operational
    (\cref{prop:match-equations,prop:premised-weaker}); pruning
    rigidly clashing branches is sound and prunes only semantically dead
    equations (\cref{prop:pruning});
  \item per-occurrence inlining of the forded inversion disjunction ---
    the shallow ``expansion'' of classified occurrences --- is sound in
    both polarities and interderivable with the axiomatic form
    (\cref{thm:expansion,prop:expansion-derivable});
  \item index uniqueness ($\Ihat{I}(w,\tup{z}) \wedge \Ihat{I}(w,\tup{z}')
    \to \tup{z} \feq \tup{z}'$) is a first-order consequence of the
    emitted axioms, while injectivity of the type former itself is
    neither expressible in the signature nor emitted --- as it must not
    be, being anti-classical (\cref{prop:index-unique},
    \cref{rem:hur});
  \item the unconditional singleton collapse is the exact point where the
    model's proof irrelevance is spent as uniqueness of identity proofs;
    Werner's guarded singleton reduction is the premise-carrying
    alternative, and both variants are sound (\cref{fact:uip},
    \cref{prop:guarded-scase}).
  \end{itemize}
\item We re-verify the shallowness checklist of the companion paper on
  the extended translation and extend it with an indexed-specific item:
  index constraints arrive as first-order equations computed at
  translation time --- never as unification problems, never as type-term
  encodings, never through an inhabitation predicate
  (\cref{sec:shallow}).
\end{enumerate}

\subsection{What soundness means here}
\label{sec:intro-soundness}

As in the companion paper, soundness is model-theoretic: if the emitted
theory first-order proves the translated goal, the goal is true in a fixed
proof-irrelevant, classical semantics of the source development
(\cref{thm:soundness}).  This is the property with operational
consequences for a hammer --- no refutable goal acquires a phantom proof,
premise-selection learning is not poisoned, and the emitted theory is
consistent --- while proof reconstruction inside the proof assistant
remains the final arbiter of each individual proof.  The semantics
validates excluded middle and proof irrelevance by construction (the ATPs
are classical; erasure forgets proofs); consequently it also validates
uniqueness of identity proofs, and \cref{sec:formulas-uip} makes precise
where the translation spends that strength and what the corresponding
reconstruction obligations are.  The converse property, completeness, is
deliberately traded away; \cref{sec:limitations} delimits what is lost.

\paragraph{The end-to-end guarantee, stated once.}
Soundness begins at $\CICi$.  The front end that carves the fragment
out of full CIC --- stratification, monomorphisation, elaboration of
native dependent matches into forded ones, dropping of proof-sorted
index positions (\cref{sec:calculus-cic}) --- is part of the pipeline,
not of the theorems, and the forded and native case styles are
intertranslatable only in CIC extended with UIP \cite{GoguenEtAl2006}.
The operational reading above is therefore relative to $\CICi$:
\cref{cor:no-refutable} says that no $\CICi$-refutable goal becomes
provable.  Composed with the front end, it reads: no goal refutable in
the user's development acquires a phantom proof, \emph{provided} the
front end is meaning-preserving --- in particular, reflects
CIC-refutability into $\CICi$-refutability --- an assumption that
silently absorbs UIP on the source side, matching the model's
validation of UIP on the target side (\cref{sec:formulas-uip}); and
proof reconstruction inside the proof assistant remains the final
arbiter of every individual proof regardless.  Neither the front end
nor its meaning preservation is formalised in this paper.  This
paragraph is the one place where the composed, conditional claim is
stated; every theorem in the body is unconditional and starts at
$\CICi$.

\subsection{Relation to the implemented translation}
\label{sec:intro-impl}

$\CICi$ and its translation are an idealisation of a design for
CoqHammer's translation pipeline, not a description of its current code.
The current implementation refuses to classify indexed families for
erasure-based treatment; declaration-level inversion axioms for inductive
predicates already carry index equations, but they are mediated by an
inhabitation atom rather than expanded per occurrence; a user-written
match on an equality proof receives no definitional axiom at all (only a
fixed list of transport constants is special-cased); and per-constructor
case equations for indexed matches are emitted without index premises,
including for statically impossible branches.  Every construction of this
paper has a direct implementation site in that pipeline: the forded
constructor telescopes of \cref{sec:translation-structural} correspond to
classifying the forded form of a family; the expansion theorem of
\cref{sec:formulas-expansion} to per-occurrence guard expansion; the
singleton elimination of \cref{sec:calculus} to the erasure of matches on
equality-like proofs; and the propositions of
\cref{sec:formulas-equations} to the choice points (index premises,
branch pruning) that the implementation exposes as options.  We indicate
the correspondences where relevant but keep the development
self-contained.
